\chapter{Conclusiones y trabajo futuro}
% TODO: Hacer. BDD de virus, integrar el proceso de @firma en la página web en vez del upload. Integración con el resto de servicios externos. 

% Este proyecto ha permitido desarrollar una plataforma de gestión de infraestructura que centraliza el control sobre usuarios, máquinas y repositorios, proporcionando una solución integral para entornos que requieren un alto grado de organización y seguridad. A lo largo del desarrollo se han aplicado buenas prácticas consolidadas en la industria, como el patrón MVC en el backend, una arquitectura basada en componentes en el frontend, diseño RESTful de la API, autenticación mediante tokens JWT y el uso de prepared statements para garantizar la seguridad frente a inyecciones SQL. La base de datos actúa como fuente de verdad única del sistema, asegurando la consistencia de la información en todo momento.
% A nivel personal, este proyecto ha puesto de manifiesto la importancia de contar con sistemas aislados, fuentes únicas de verdad y un control riguroso sobre las máquinas, con el fin de minimizar riesgos y evitar la pérdida de datos. El desarrollo por incrementos ha permitido iterar sobre las funcionalidades de forma ordenada, validando cada avance antes de abordar el siguiente.
%
%
% Aunque el proyecto cubre las funcionalidades esenciales alcanzado un producto mínimo viable, existen diversas líneas de mejora y expansión que quedan abiertas para iteraciones futuras.
%
% En el ámbito de la integración con GitLab, se contempla ampliar las funcionalidades al nivel de grupos, permitiendo una gestión más completa de los repositorios y sus miembros. En relación a los codeservers, se prevé la creación de colas de ejecución que permitan gestionar y ordenar los trabajos de manera eficiente, evitando conflictos por uso simultáneo de recursos. Además, se plantea el autodespliegue de los servicios de la página web, reduciendo la intervención manual en los procesos de puesta en producción. A nivel de arquitectura de la API, se tiene como objetivo alcanzar el máximo nivel del Richardson Maturity Model, incorporando controles HATEOAS para lograr una API completamente autodescriptible y navegable. Por último, se contempla la integración con una base de datos de firma de virus, con el fin de reforzar la seguridad del sistema frente a software malicioso en las peticiones, las upload del usuario. Al igual que se podría añadir sistema de autenticación de doble factor al sistema dando una capa de seguridad adicional.
%
% Además de lo anterior, quedaría añadir las integraciones externas planteadas de manera teórica. Al igual que añadir las mejoras planteadas en los resultados de las pruebas de usabilidad.
%

Este proyecto ha permitido desarrollar una plataforma de gestión de infraestructura que centraliza el control sobre usuarios, máquinas y repositorios, proporcionando una solución integral para entornos que requieren un alto grado de organización y seguridad. A lo largo del desarrollo se han aplicado buenas prácticas consolidadas en la industria, entre las que destacan el patrón MVC en el backend, una arquitectura basada en componentes en el frontend, el diseño RESTful de la API, la autenticación mediante tokens JWT y el uso de prepared statements para garantizar la seguridad frente a inyecciones SQL. La base de datos actúa como única fuente de verdad del sistema, asegurando la consistencia de la información en todo momento.
A nivel personal, este proyecto ha puesto de manifiesto la importancia de contar con sistemas aislados, fuentes únicas de verdad y un control riguroso sobre las máquinas, con el fin de minimizar riesgos y evitar la pérdida de datos. El desarrollo por incrementos ha permitido iterar sobre las funcionalidades de forma ordenada, validando cada avance antes de abordar el siguiente. Aunque el proyecto cubre las funcionalidades esenciales alcanzando un producto mínimo viable, existen diversas líneas de mejora y expansión que quedan abiertas para iteraciones futuras.

Por otro lado, se destaca la relevancia de integrar un sistema de peticiones centralizado para el acceso a los recursos del laboratorio, lo que garantiza una trazabilidad completa. Bajo este enfoque, cada servicio desplegado queda vinculado inequívocamente a su solicitud correspondiente, optimizando el control y el seguimiento de su ciclo de vida.


En el ámbito de la integración con GitLab, se contempla ampliar las funcionalidades al nivel de grupos, lo que permitiría una gestión más completa de los repositorios y sus miembros. En relación con los code-servers, se prevé la incorporación de colas de ejecución que permitan gestionar y ordenar los trabajos de manera eficiente, evitando conflictos derivados del uso simultáneo de recursos. Asimismo, se plantea el autodespliegue de los servicios de la aplicación web, con el objetivo de reducir la intervención manual en los procesos de puesta en producción.
Desde el punto de vista de la arquitectura de la API, se tiene como objetivo alcanzar el nivel máximo del Richardson Maturity Model mediante la incorporación de controles HATEOAS, lo que permitiría disponer de una API completamente autodescriptible y navegable. En materia de seguridad, se contempla la integración con una base de datos de firmas de virus para reforzar la protección del sistema frente a software malicioso en las peticiones y en los archivos subidos por los usuarios, así como la incorporación de un sistema de autenticación de doble factor que añadiría una capa de seguridad adicional.

Por último, quedan pendientes de implementación las integraciones externas planteadas en la fase de diseño, así como las mejoras identificadas a partir de los resultados obtenidos en las pruebas de usabilidad.
